% STATUS: written 2026-11-09; verified against the full text 2026-11-09 (issue #1325).
% Sources read in full, not at the abstract:
%   Benkovic 2011, Theorem 2.4 (p. 823) and the abstract (p. 819);
%   Lu 2007, Theorem 1.1 (p. 882);
%   Li and Fang 2013, the abstract (p. 653);
%   Ji and Wang 2005, the abstract (p. 399) and section 1 (p. 399).
% Correction 2026-11-09: the TUHF paper was filed in the bibliography as
% "Zhu, Jun, Proc. Amer. Math. Soc. 133 (2005)". The paper is by Ji and Wang,
% Linear Algebra Appl. 403 (2005) 399--408, doi 10.1016/j.laa.2005.02.004, and the
% key is now JiWang2005TUHF.
% Correction 2026-11-09: Benkovic Theorem 2.4 has three terms, not two. The
% earlier text said "a standard Lie derivation plus an antiderivation" and
% dropped the central term.
\section{Lie triple derivations across the classes of algebras}\label{sec:lie-type:triple-across-classes}

A Lie triple derivation satisfies the triple rule and need not be a Lie
derivation, because the two differ by a map into the centre that vanishes on
brackets.
Section~\ref{sec:lie-type:triple-and-n} stated the rule and gave the
decomposition on von Neumann algebras.
This section carries the same question to the concrete classes, where the answer
has one shape: an inner or standard part plus a central defect. What changes
between the classes is which derivation counts as standard and whether the defect
can be written down.

Take the algebra $T_{n}(C)$ of upper triangular matrices over a commutative
unital ring, with a $2$-torsion free unital bimodule $M$.
Every Lie triple derivation $d$ of $T_{n}(C)$ into $M$ is a sum of three terms,
\begin{equation*}
  d = \Delta + \delta + \tau .
\end{equation*}
Here $\Delta$ is a derivation and $\delta$ is an antiderivation that vanishes on
the diagonal matrices. The third term $\tau$ carries $T_{n}(C)$ into the centre
of $M$ and sends every commutator to zero
\cite[Theorem~2.4, p.~823]{Benkovic2011TripleTriangular}.
The three terms are visible separately in the abstract, which names the standard
Lie derivation and the antiderivation and leaves the central term to the theorem.
The antiderivation appears because the algebra is triangular.

The nest algebra case is known in both directions. Take a nest $N$ of projections
on a complex Hilbert space and a weak$^{\ast}$ closed algebra $B$ containing the
nest algebra $T(N)$.
A linear map $\delta$ of $T(N)$ into $B$ is a Lie triple derivation if and only if
\begin{equation*}
  \delta(A) = [S,A] + f(A)I \qquad \text{for all } A \in T(N),
\end{equation*}
where $S \in B$ and $f$ is a linear functional that vanishes on every second
commutator $[[X,Y],Z]$ \cite[Theorem~1.1, p.~882]{Lu2007LieTripleNest}.
The if-and-only-if is the content: the inner part plus that functional is the
whole class rather than one description among several.

The same algebras were treated without linearity. Write $\operatorname{Alg}N$ for
the nest algebra of a nest on a Banach space, and suppose that a nontrivial
element of the nest is complemented. Let $\Phi$ be a Lie triple derivable map of
$\operatorname{Alg}N$ into $B(X)$. Then for all $A,B$ in the algebra there is a
scalar $\lambda_{A,B}$ with
$\Phi(A+B)=\Phi(A)+\Phi(B)+\lambda_{A,B}I$. Moreover $\Phi = D + \tau$, where $D$
is an additive derivation. The map $\tau$ carries the algebra into the scalar
multiples of the identity, has the same additive defect, and vanishes on second
commutators \cite[Abstract, p.~653]{LiFang2013LieTripleNest}.
On the same algebra a Jordan derivable map is an additive derivation, with no
defect at all \cite[Abstract, p.~653]{LiFang2013LieTripleNest}.
The two derivable conditions therefore differ on nest algebras. On rings the
Jordan condition is the weaker one, so this is the case worth naming.

A TUHF algebra is a triangular algebra of Hilbert-space operators. For a
continuous Lie triple derivation $L$ of one there is an associative derivation $D$
with $L = D + \lambda$.
The map $\lambda$ carries the algebra into its centre and sends every bracket of
operators to zero \cite[Abstract, p.~399]{JiWang2005TUHF}.
That is the triangular shape again, restricted to the continuous maps.

Full matrix algebras give the same shape without linearity. On $M_{n}(R)$ over a
$2$-torsion free commutative ring with identity, where $n\ge 2$, every nonlinear
Lie triple derivation is a sum of three terms: an inner derivation, an additive
induced derivation, and a functional that annihilates all second commutators
\cite[Abstract, p.~77]{XiaoWei2013NonlinearLieTriple}.

A ring with a nontrivial idempotent gives the same decomposition, linearity apart.
Under a mild condition, a Lie triple derivable map $\Phi$ of such a ring into
itself satisfies $\Phi(A+B)=\Phi(A)+\Phi(B)+z_{A,B}$ for a central element
$z_{A,B}$ depending on the pair. On a prime ring of characteristic not two,
$\Phi = D + \tau$, where $D$ is an additive derivation into the central closure
and $\tau$ carries the ring into the extended centroid with the same additive
defect \cite[Abstract]{LiFangLuWang2014LieTripleDerivable}.

Three of the four results say the same thing. Each is a standard derivation
together with a map into the centre that vanishes on brackets.
The fourth says the two derivable conditions differ.
The section exists so that a reader who meets the triple rule in one class can
see what it becomes in the others.

On a finite von Neumann algebra with no central summand of type $I_1$, every second nonlinear mixed Lie triple derivation is an additive $*$-derivation, and each local and 2-local second Lie triple derivation is one too \cite[Abstract]{SecondNonlinearMixedLieTriple2020}.

A map of a standard operator algebra satisfying the multiplicative Lie triple rule is shown to be a Lie triple derivation plus a central term~\cite{WlMultiplicativeLieTriple2018}.

A map of a standard operator algebra satisfying the multiplicative Lie triple rule is shown to be a Lie triple derivation plus a central term~\cite{WlMultiplicativeLieTriple2018}.
